%────────────────────── Capítulo 3 — Mitigaciones ──────────────────────
% Enumeración de posibles mitigaciones y su impacto.
\chapter{Mitigaciones}
\label{cap:mitigaciones}

El ataque del \cref{cap:descripcion} necesita dos cosas: que el consumo en un
instante concreto dependa de un valor intermedio que el atacante puede predecir,
y que pueda repetir la medida las veces suficientes para que la estadística
saque esa dependencia del ruido. Las contramedidas atacan una de las dos.
Mangard, Oswald y Popp las agrupan en dos familias \cite{mangard2007power}: el
\emph{ocultamiento} (\emph{hiding}), que reduce la relación señal-ruido haciendo
el consumo constante o impredecible en el tiempo, y el \emph{enmascaramiento}
(\emph{masking}), que hace que los valores que se procesan dejen de ser
predecibles aunque el consumo siga dependiendo de ellos. A estas se suman las
medidas de protocolo, como la renovación de claves, que limitan cuántas trazas
puede reunir el atacante con una misma clave. La revisión de Mishra y Sahay
recoge las tres \cite[secs.~II.C.5 y III.D]{mishra2025modern}.

Ninguna de ellas es una garantía absoluta. Cada una sube el coste del ataque
---más trazas, mejor equipo, un modelo más elaborado---, y casi todas tienen un
coste para el defensor en tiempo de ejecución, área, energía o complejidad.
Las secciones siguientes describen cada una, su coste y sus límites, y qué
efecto tendría sobre el ataque de este trabajo.

\section{Enmascaramiento (\emph{masking})}
\label{sec:masking}

El enmascaramiento combina cada valor sensible con un número aleatorio nuevo en
cada ejecución, la \emph{máscara}, y opera sobre el valor enmascarado. En el
enmascaramiento booleano, en lugar de $v$ se manipula $v \oplus m$, con $m$
uniforme y desconocida para el atacante. Cada valor que pasa por el procesador
es, por separado, independiente de la clave, y la predicción
$\operatorname{HW}(\mathrm{SBox}(p_j \oplus k))$ del CPA deja de correlar con
el consumo de cualquier instante. Chari et al.\ formalizaron la idea dividiendo
cada valor en varias partes aleatorias y demostraron que, si cada parte se mide
con ruido, el número de trazas necesario crece exponencialmente con el número de
partes \cite{chari1999masking}.

En AES, las operaciones lineales (\texttt{AddRoundKey}, \texttt{ShiftRows},
\texttt{MixColumns}) se aplican a cada parte por separado. La S-box no es
lineal y obliga a algo más: en software, lo habitual es calcular en cada
ejecución una tabla enmascarada $\mathrm{SBox}'(x \oplus m) = \mathrm{SBox}(x)
\oplus m'$ \cite{herbst2006aes}. Eso cuesta memoria para la
tabla, tiempo para recalcularla y una fuente de números aleatorios en el
dispositivo.

El límite principal es que el enmascaramiento de primer orden solo protege
frente a ataques que miran un instante. Un atacante que combina las muestras de
dos instantes ---el de la máscara y el del valor enmascarado--- recupera una
dependencia con la clave: es el DPA de segundo orden, que Messerges ya aplicó en
2000 contra software enmascarado \cite{messerges2000second}. Más partes elevan
el orden del ataque necesario, y con ruido suficiente su coste, pero también el
del cifrado. Además, la protección depende de que las partes no se combinen
nunca en el hardware sin querer, por ejemplo al sobrescribir un registro que
contenía una con la otra, y eso es difícil de garantizar en un procesador real.
% TODO: fuente concreta (capítulo de Mangard et al. sobre ataques a
% implementaciones enmascaradas, o artículo específico).

Contra el ataque de este trabajo, una implementación enmascarada correcta haría
fallar el CPA de primer orden tal como está escrito, independientemente del
número de trazas.

\section{Aleatorización y \emph{shuffling}}
\label{sec:shuffling}

La aleatorización no oculta el valor, sino el momento en que se procesa. El CPA
compara las trazas instante a instante, así que supone que la S-box del byte $j$
se calcula siempre en la misma muestra. Hay dos formas de romper esa
alineación en software:

\begin{itemize}
  \item \emph{Shuffling}: procesar los 16 bytes del estado en un orden
    aleatorio distinto en cada cifrado. Las operaciones de AES son las mismas
    en cualquier orden, y el coste es pequeño: generar una permutación por
    cifrado \cite{herbst2006aes}.
  \item Retardos aleatorios: insertar entre las operaciones bucles de duración
    aleatoria, de modo que todo lo que viene detrás se desplaza
    \cite{clavier2000countermeasures}. Cuesta tiempo de ejecución, en
    proporción a la longitud media de los retardos.
\end{itemize}

El efecto sobre el CPA se puede estimar con el razonamiento de la
\cref{sec:cpa}. Si la S-box de un byte cae con la misma probabilidad en
cualquiera de $m$ posiciones, en cada una solo está presente en una de cada $m$
trazas, y su correlación se reduce aproximadamente a $\rho/m$. Como el número
de trazas crece como el inverso del cuadrado de la correlación, harían falta
del orden de $m^2$ veces más. Con el \emph{shuffling} de 16 bytes, $m = 16$. El
atacante puede compensarlo en parte sumando las $m$ ventanas posibles antes de
correlar, como proponen Clavier et al.\ frente a los retardos aleatorios
\cite{clavier2000countermeasures}; si el ruido de las ventanas es independiente,
la correlación queda en torno a $\rho/\sqrt{m}$ y el factor baja a unas $m$
veces. Estas estimaciones salen de un modelo simple y no se han comprobado
sobre trazas: medirlas es el objetivo de la fase opcional de contramedidas.
% TODO: si se hace la fase 6, comparar con las trazas necesarias medidas.

El límite de la aleatorización es que no elimina la fuga, solo la diluye: con
suficientes trazas, el ataque vuelve a funcionar. Por eso se combina con el
enmascaramiento, como en la implementación de Herbst et al.\ para tarjetas
inteligentes de 8 bits \cite{herbst2006aes}. Además, si el SPA permite
reconocer los retardos en una traza individual, se pueden eliminar alineando las
trazas antes del ataque.

\section{Lógica balanceada}
\label{sec:logica-balanceada}

Las lógicas balanceadas atacan la causa física de la fuga. En la lógica de
doble raíl con precarga, cada bit se representa con dos señales
complementarias que se precargan a un valor fijo en cada ciclo y después se
evalúan: en cada ciclo cambia exactamente una de las dos, sea cual sea el
dato, y el consumo es en principio independiente del valor
\cite{tiri2004wddl}. WDDL (\emph{wave dynamic differential logic}) implementa
esa idea con celdas estándar, de modo que se puede aplicar en ASIC o en FPGA
\cite{tiri2004wddl}; Mishra y Sahay la citan entre las contramedidas hardware
frente al DPA \cite[sec.~III.D]{mishra2025modern}.

El coste es alto: se duplican las señales y se añade la fase de precarga, lo que
aumenta el área, el consumo y el tiempo de cada operación. El límite es que el
equilibrio nunca es perfecto, porque las dos líneas de cada par no tienen
exactamente la misma capacidad tras el rutado, y esa diferencia residual
vuelve a depender del dato \cite{mangard2007power}.
% TODO: fuente: precisar el capítulo de Mangard et al. (ocultamiento).

Esta contramedida se decide al diseñar el circuito. En un microcontrolador
comercial como el ATmega328P, que no la incorpora, no se puede aplicar desde el
firmware.

\section{Inyección de ruido}
\label{sec:ruido}

Otra forma de ocultamiento es añadir consumo que no depende de los datos: un
generador de ruido en la alimentación, circuitos que conmutan al azar o, en
software, operaciones falsas intercaladas con las reales
\cite[sec.~III.D]{mishra2025modern}. El objetivo es bajar la relación
señal-ruido de la medida.

Es una de las medidas más baratas, pero su efecto sobre el CPA es limitado
porque el ruido independiente de los datos se cancela al promediar. La
simulación de la \cref{sec:validacion-simulada} lo ilustra: al pasar de una
relación señal-ruido de 0,2 a 0,02, las trazas que necesita el 90~\% de los
ensayos pasan de 295 a 2220, unas 7,5 veces más (\cref{tab:sim-ruido}). El ataque se encarece de forma lineal, no exponencial
como con el enmascaramiento, y un atacante con tiempo y un buen equipo de
captura lo compensa. Por eso el ruido se usa como complemento de otras medidas y
no como protección única.

\section{Renovación de claves (\emph{re-keying})}
\label{sec:rekeying}

La renovación de claves no toca la implementación del cifrado, sino el uso que
se hace de él. Si cada clave cifra menos bloques de los que el atacante
necesita para recuperarla, el CPA no llega a converger. En el esquema de
\emph{fresh re-keying} de Medwed et al., cada sesión cifra con una clave nueva
que se deriva de la clave maestra y de un número aleatorio público; la función
de derivación se elige para que sea fácil de proteger, de modo que solo ella
necesita contramedidas costosas, y el cifrado de cada sesión se puede dejar sin
proteger \cite{medwed2010fresh}. Mishra y Sahay mencionan también la opción de
fijar el intervalo de renovación a partir del número de trazas que necesita un
ataque \cite[sec.~II.C.5]{mishra2025modern}.

Su límite es que desplaza el problema en lugar de eliminarlo: la clave maestra
sigue en el dispositivo y la derivación sigue siendo un objetivo. Además,
requiere cambiar el protocolo, porque las dos partes tienen que sincronizar la
clave de cada sesión, y no protege frente a ataques que necesitan muy pocas
trazas, como los de plantillas \cite[sec.~II.B.4]{mishra2025modern}.

En el modelo simulado, con una relación señal-ruido de 0,1, la mediana de
trazas para recuperar la clave completa está entre 296 y 413
(\cref{tab:sim-ruido}). Una clave que se renovara cada pocas decenas de
cifrados quedaría fuera del alcance de ese ataque, pero el umbral real depende
del montaje y no se conoce hasta medir.

\section{Tiempo constante no es protección frente al consumo}
\label{sec:tiempo-constante}

Una implementación de tiempo constante ejecuta las mismas instrucciones, en el
mismo orden y con la misma duración, sea cual sea la clave. Protege frente a los
ataques de tiempo \cite{kocher1996timing} y frente al SPA basado en el flujo de
ejecución, pero no frente al DPA ni al CPA: el consumo de cada instrucción sigue
dependiendo de los valores que manipula. De hecho, un código de tiempo constante
favorece al atacante de consumo, porque produce trazas perfectamente alineadas.
Es exactamente la situación del modelo simulado de la
\cref{sec:validacion-simulada}, en el que el CPA recupera la clave.

\section{Resumen}
\label{sec:resumen-mitigaciones}

La \cref{tab:mitigaciones} resume las contramedidas. Su coste y su eficacia
se indican de forma cualitativa, porque dependen de la implementación y no se
han medido en este trabajo.

\begin{table}[htbp]
  \centering
  \caption{Resumen de las contramedidas frente al análisis de consumo. Las
    valoraciones son cualitativas y proceden de las referencias citadas en
    cada sección.}
  \label{tab:mitigaciones}
  \small
  \begin{tabularx}{\textwidth}{>{\raggedright\arraybackslash}p{2.9cm}
                               >{\raggedright\arraybackslash}p{1.8cm}
                               >{\raggedright\arraybackslash}X
                               >{\raggedright\arraybackslash}X}
    \toprule
    Contramedida & Nivel & Efecto sobre el CPA & Límite principal \\
    \midrule
    Enmascaramiento & Algoritmo & Anula el CPA de primer orden &
      Ataques de orden superior; combinación accidental de las partes \\
    \emph{Shuffling} y retardos & Algoritmo & Diluye el pico; más trazas &
      No elimina la fuga; integración de ventanas y realineado \\
    Lógica balanceada & Circuito & Reduce la fuga en origen &
      Desequilibrio residual; solo en diseño de hardware \\
    Inyección de ruido & Circuito o software & Baja la SNR; más trazas &
      El ruido se cancela al promediar \\
    Renovación de claves & Protocolo & Limita las trazas por clave &
      La derivación de claves sigue expuesta; cambia el protocolo \\
    \bottomrule
  \end{tabularx}
\end{table}
